\chapter{The Protocol Is the Program}
\label{ch:language}
\label{ch:vision}

\index{LeanFM language}
\index{protocol}
LeanFM is a programming language whose programs are protocols. A protocol names
the roles in a conversation, the messages they exchange, the order those messages
come in, who decides at each choice, what is assumed about decisions nobody else
can see, and what the system promises. Everything else in a LeanFM project is
generated from the protocol or checked against it: the diagrams people read, the
requirements and implementation mapping reviewed in git, the OpenAPI and protobuf
interfaces, the typestate code each role is written on, and the verdicts on what
the running system logs.

This book documents that language: how to write it, and what it means.

\section{Who writes it}

Most LeanFM is written by Claude. A person says what the software must do, in
ordinary words. Claude writes the protocol, LeanFM checks it, and the person
looks at what comes back: a sequence diagram, each role's obligations, a
counterexample. They push back, and the protocol changes. Code comes last,
written in one pass from the checked protocol.

Claude learns the language from LeanFM itself. The LeanFM server speaks the Model
Context Protocol, so an agent that connects to it is handed the language
reference and the tools that check, draw, simulate and judge a protocol. MCP is
the language's tutorial. Appendix~\ref{ch:mcp-tools} lists the tools.

\section{A protocol}

\index{PlaceOrder}
Here is a complete program: a client places an order, the shop quotes or says it
is sold out, the client pays, and the bank charges the card or declines it. The
structures at the top are its message types.

\begin{lstlisting}[language={},basicstyle=\ttfamily\footnotesize,breaklines=true]
/-- Placing an order: quote, pay, charge the card. -/
protocol PlaceOrder (Client, Shop, Bank) {
  o: CreateOrder from Client to Shop route "POST /orders";
  choice {
    q: Quote from Shop to Client cause o status 201;
    p: Pay from Client to Shop cause q where p.amount = q.price;
    c: Charge from Shop to Bank cause p;
    choice {
      d: Charged from Bank to Shop cause c;
      Receipt from Shop to Client cause d status 200;
    } or {
      x: Declined from Bank to Shop cause c;
      PaymentDeclined from Shop to Client cause x status 402;
    }
  } or {
    SoldOut from Shop to Client cause o status 409;
  }
}
\end{lstlisting}


LeanFM draws it. This diagram is LeanFM's own rendering of the text above, not a
drawing of it; when the text changes, so does the picture.

\begin{figure}[htbp]
\centering
\AccessibleFigure{Sequence diagram of PlaceOrder. Client sends CreateOrder to Shop on POST /orders. In an alt frame the Shop decides. In one branch Shop sends Quote 201 to Client. Client sends Pay to Shop when the amount paid equals the price quoted. Shop sends Charge to Bank. In an inner alt frame the Bank decides: Charged and then Receipt 200 to Client. Or Declined and then PaymentDeclined 402 to Client. In the other branch Shop sends SoldOut 409 to Client. Responses are dashed. Each call is an activation bar on the callee.}{%
\ThemedDiagram{width=.82\textwidth}{lang-placeorder}}
\caption{\texttt{PlaceOrder} as LeanFM draws it. Each \texttt{alt} frame names the
role that decides; responses are dashed, and each call is an activation bar on
its callee.}
\label{fig:placeorder}
\end{figure}

Each line is one message: a label, a sender and a receiver. The rest is optional.
\texttt{o:} names the message so later ones can point at it. \texttt{cause o} says
this message answers \texttt{o}. \texttt{route} and \texttt{status} make the
exchange an HTTP call and its response. \texttt{where} states a condition on the
payloads, and conformance runs it against what the running system logs
(Chapter~\ref{ch:execution}): a client that pays 12 against a quote of 10 does not
conform.

\section{Message types}

\index{message type}
\index{structure}
Every message has a type, written as a Lean structure in the session type. A
message carries the structure its label names: \texttt{Pay} carries an
\texttt{amount} and a \texttt{card\_secr}. A label that names no structure, such as
\texttt{SoldOut}, carries nothing, and so does an event. A field is a
\texttt{Nat}, \texttt{Int}, \texttt{String}, \texttt{Bool}, another structure, or a
\texttt{List} or \texttt{Option} of one of those; \texttt{x y : Nat} declares two
fields, as in Lean, so a move in a game is \texttt{structure Play where x y : Nat}
and a pass is \texttt{structure Pass}.

The session type is the source, and the rest is written from it, the way an
OpenAPI file is written for an HTTP interface: the \texttt{.proto}, the Lean module
(Chapter~\ref{ch:generated}) and the OpenAPI document. Implementations use those
directly. Fields are numbered in the order written, and that number is the
field's protobuf number, so new fields go at the end:

\begin{lstlisting}[language={},basicstyle=\ttfamily\footnotesize,breaklines=true]
// What the client orders.
message CreateOrder {
  string item = 1;
  uint64 quantity = 2;
}

// The shop's price for the order, in cents.
message Quote {
  uint64 price = 1;
}

// The client pays; the card number is a secret, never logged.
message Pay {
  uint64 amount = 1;
  string card_secr = 2;
}

message Charge {
  uint64 amount = 1;
  string card_secr = 2;
}
\end{lstlisting}


The same structures give each actor's HTTP interface, written as OpenAPI beside
the \texttt{.proto} (Chapter~\ref{ch:generated} shows the shop's).

An existing \texttt{.proto} can stand in for the structures. LeanFM reads only the
part with an exact Lean type: unsigned integers are \texttt{Nat}, signed ones
\texttt{Int}, and every field keeps its number. \texttt{float} and
\texttt{double} are refused because they are not exact (a price is an integer
number of cents), and so are \texttt{import} (LeanFM reads only the text it is
given), nested messages, \texttt{enum}, \texttt{oneof} and \texttt{map}.

The types are what the conditions and the logs are checked against. A
\texttt{where} condition is typed when the text is read, before anything runs:
\texttt{p.amount = q.price} compares two declared \texttt{Nat} fields;
\texttt{p.total = q.price} is refused because \texttt{Pay} has no
\texttt{total}, and \texttt{p.amount = "ten"} because a number is not a string. A
name ending in \texttt{\_secr} is a secret, and no condition may read one,
because the fields conditions read are logged. A log is checked against the same
types: a field the message does not declare, a value of another type, or a secret
in a log is refused.

\section{Statements}

\index{statements}
A protocol is a list of statements. These are all of them.

\begin{center}
\small
\begin{tabular}{@{}p{.43\textwidth}p{.5\textwidth}@{}}
\toprule
Statement & Meaning \\
\midrule
\texttt{x: L from A to B cause y route "\ldots" status 201 where e;} & a message; everything after \texttt{to B} is optional \\
\texttt{L at A;} & something role \texttt{A} does on its own, such as a timer firing \\
\texttt{choice \{ \ldots \} or \{ \ldots \}} & exactly one branch happens \\
\texttt{choice at A using declared("N") \{ probability p: \ldots \}} & the same, with \texttt{A}'s decision assumed to have these probabilities (Chapter~\ref{ch:assumptions}) \\
\texttt{par \{ \ldots \} and \{ \ldots \}} & both branches run; what follows waits for all of them \\
\texttt{join 2 of \{ \ldots \} and \{ \ldots \} and \{ \ldots \}} & what follows waits for any two \\
\texttt{rec Loop \{ \ldots\ continue Loop; \}} & a loop; \texttt{continue} is the last statement of its block \\
\texttt{sessions \{ \ldots \}} & any number of independent instances, as a server handles requests \\
\texttt{include P A B C;} & an earlier protocol, its roles renamed in order \\
\texttt{require per q \{ \ldots \}} & what the protocol promises per occurrence of message \texttt{q} (Chapter~\ref{ch:assumptions}) \\
\bottomrule
\end{tabular}
\end{center}

Statements after a choice continue every branch. Names are Lean identifiers, and
\texttt{--} and \texttt{/- -/} comments are allowed. Several protocols may share a
text, so one can include another; the last is the one checked. This is how a large
protocol is broken up: each part is its own named protocol, written as if it were
the whole conversation, and the top-level protocol assembles them with
\texttt{include}, renaming roles as it goes. Each part is checked on its own, and the
top-level diagram draws each as a frame that names it. A protocol may
carry amounts, delays and rates on its messages; those annotations never change
its structure, and Chapter~\ref{ch:assumptions} covers them.

\index{choice!decider}
A plain \texttt{choice} does not name who decides, because only one answer can
be right: the role that sends the first message of every branch. In
\texttt{PlaceOrder} the shop decides whether to quote, and the bank whether to
charge. If two roles could each start a branch, both could act at once, and
LeanFM reports the race (Chapter~\ref{ch:meaning}).

\subsection*{A session type is a Lean file}

\index{Lean!session-type file}
A \texttt{.leanfm} file is Lean. Its first line, \texttt{import
LeanFM.Session.Syntax}, loads the library that defines the \texttt{protocol}
command and the \texttt{require} command, and from the LeanFM directory

\begin{lstlisting}[language={},basicstyle=\ttfamily\small]
lake env lean examples/sessions/PlaceOrder.leanfm
\end{lstlisting}

\noindent elaborates it like any other Lean file. Every construct in this chapter
and in Chapter~\ref{ch:assumptions} parses: amounts, delays, clocks, obligations,
accruals, declared choices, bounded loops, keyed sessions and the requirements
after a protocol. Each label is a Lean type. A label that no structure declares,
such as \texttt{SoldOut}, is a message that carries nothing, and the command
declares the empty structure for it. Each \texttt{where} condition is
type-checked as a Lean proposition over the structures, so \texttt{p.total} is a
Lean error in the editor as well as a refusal from the tools. The condition stays
in LeanFM's own small language of comparisons, arithmetic and \texttt{and},
\texttt{or}, \texttt{not}: a function call is not part of it and does not parse.
Because the file is Lean, a protocol, role, label or message id may not be a
Lean keyword, and a protocol may not share its name with one of its labels.

There is one grammar, written with Lean's own parser, and no second parser. A
Lean file is parsed with it, and the tools parse the text a client sends with the
same parsers, by Lean's token rules, without elaborating anything: nothing a client
sends is compiled or run in the server. One reading then turns the parsed text
into statements, with the same checks and limits, whichever way the text arrived.
The tools type the \texttt{where} conditions themselves and run them with
LeanFM's interpreter, so text from any client cannot execute code, and nesting
deep enough to exhaust the parser is refused before parsing starts. The test
suite elaborates every session-type file in the repository, and the text the tools
write back, and compares the protocol each defines with the one the tools read:
the same roles, diagram, issues and local types. Anyone who wants more isolation
than that can run LeanFM itself in a container or a virtual machine.

\section{Only what crosses a boundary}

\index{hidden state}
A protocol says nothing about what happens inside a role. Whether the shop keeps
its stock in a database, a cache or a spreadsheet is left free by not being
written down. This is deliberate. A requirement that depends on a role's internal
state is a design decision that leaked out, and checking it costs state space:
every variable the checker must track multiplies the states it explores. A
TLA+ module of the same order queue with a few cache bits and a clock that cannot
affect the order shows the cost (counts from TLC, via \texttt{make tla-check}):

\begin{center}
% Generated from TLC logs by scripts/check_tla_example.py.
\begin{tabular}{@{}lrrr@{}}\toprule
Model & Cache bits & Clock bound & Distinct states \\\midrule
Baseline & 0 & 0 & 26 \\
Cache4 & 4 & 0 & 416 \\
Cache8 & 8 & 0 & 6,656 \\
Clock7 & 0 & 7 & 208 \\
\bottomrule\end{tabular}

\end{center}

Each irrelevant bit doubles the states, and an unbounded clock makes them
infinite \citep{lamport2002specifying}. TLA+ can omit such state too; LeanFM
starts from the conversation, so there is nowhere to put it. The state that
matters lives in the messages, and an execution is the stream of them
(Chapter~\ref{ch:execution}).

\section{Where the language comes from}

\begin{enumerate}[label=\arabic*.]
  \item \textbf{Session types} give the backbone: one global description of a
        conversation, projected to each role's obligations, with conditions that
        rule out choices nobody can implement
        \citep{honda1993types,honda2008multiparty,yoshida2013scribble}.
  \item \textbf{Language theory} gives the reading of a protocol as a grammar
        whose terminals are messages, so the legal conversations are its
        language. Dana Richards and Henry Hamburger, the author's college
        professors, are cited in \citet{richards2017logic} as where the author
        learned these methods; for symbolization, validity and modal reasoning
        we also draw on \citet{gensler2017introduction}.
  \item \textbf{Protocol analysis} contributes explicit attackers, adversarial
        choices and trace semantics. The CSP treatment of
        \citet{ryan2001modelling}, including Steve Schneider, is where the author
        learned it.
  \item \textbf{Temporal logic} turns expectations into questions about every
        possible future, or some \citep{clarke1986automatic}.
\end{enumerate}

LeanFM adds two things the classic theory leaves out. Decisions nobody else can
see carry stated probabilities, and the protocol states what it promises, as
minimums and maximums judged over an observation period. Both are
Chapter~\ref{ch:assumptions}.
